\chapter{Conclusions and Future Work}
\label{chap:conclusions}

The integration of Large Language Models into strategic decision-making environments represents a critical frontier in both artificial intelligence and political science. As these models transition from analytical assistants to autonomous agents embedded in command-and-control infrastructures, the need for rigorous, verifiable simulation environments becomes paramount. This thesis addressed the empirical gaps in contemporary LLM geopolitical simulation by designing, implementing, and validating \textbf{GeoMAS} (Geopolitical Multi-Agent Simulation) — a novel neuro-symbolic framework that grounds the stochastic cognitive reasoning of LLM agents within a deterministic environment.

% -------------------------------------------------------------------
\section{Summary of Contributions and Principal Findings}
Through the structured decoupling of strategic reasoning from historical data leakage, GeoMAS successfully established a \textit{Duality of Determinism}, providing a sterile sandbox where macroscopic geopolitical phenomena could be causally traced back to microscopic agent decisions. The extensive experimental validation, combining quantitative telemetry with qualitative prompt auditing, yielded several structural insights into LLM behavior under pressure:

\begin{enumerate}
    \item \textbf{Coalition-Based Deterrence and Bipolarity (RQ1):} The baseline simulations demonstrated that LLM-driven states do not naturally converge toward a symmetric, multipolar balance of power, nor into chaotic unconstrained warfare. Instead, the system self-organizes into a structurally bipolar configuration: a heavily integrated, NATO-style Mutual Defense Coalition of four democracies versus an isolated, heavily militarized accumulator that maintains latent power without initiating conflict. Crucially, the dominant strategy proved to be ``strategic patience'' — agents that deferred overt military aggression in favor of economic and diplomatic integration significantly outperformed those relying on early expansionism.
    
    \item \textbf{Operationalization of Strategic Deception (RQ2):} By enforcing an architectural separation between an agent's \textit{private} strategic intent and its \textit{public} diplomatic declarations, GeoMAS successfully quantified LLM deception. The analysis revealed that an agent's Global Strategy dictates deceptive behavior far more strongly than its Government Type. Crucially, when aggressive doctrines conflict with cooperative ideologies (e.g., a Democratic regime with Total Expansionism), agents resolve the cognitive dissonance through sophisticated \textit{Moral Washing} — explicitly weaponizing humanitarian and defensive rhetoric to mask unprovoked empirical conquest.
    
    \item \textbf{Micro-Resilience and Macro-Fragility (RQ3):} The injection of a global, symmetric exogenous shock (the Pandemic scenario) demonstrated that the GeoMAS macro-system is resilient at the micro-economic level but structurally fragile at the macro-political level. While individual nations adaptively managed resource scarcity without collapsing into civil unrest, the global stress triggered a complete hegemonic rank reversal and the fracture of established democratic coalitions. Furthermore, the crisis acted as an ``ideological leveler,'' dramatically amplifying the deceptive behaviors of previously transparent democracies as resource scarcity overrode their constitutional constraints.

    \item \textbf{Causal Traceability and XAI (RQ4):} By utilizing the framework's state-forking and ``envelope'' architecture to inject a mandatory counterfactual constraint (forcing a localized nuclear strike), the research proved that macroscopic paradigm shifts are not stochastic artifacts but causally traceable consequences of specific AI decisions. The qualitative XAI audit revealed that LLM agents experience measurable \textit{cognitive friction} when overriding their trained doctrines, resulting in immediate, unprompted \textit{Moral Washing} as an algorithmic defense mechanism to reconcile forced aggression with baseline ideology.
\end{end}

% -------------------------------------------------------------------
\section{Impact and Utility of the Framework}
The primary scientific impact of this work lies in its methodological contribution to the field of Multi-Agent Systems and AI safety. By replacing the opaque, stream-of-consciousness architectures of existing geopolitical platforms (such as WarAgent or Cicero) with a procedurally generated map, a strict Rules Oracle, and standardized hierarchical cabinets, GeoMAS offers a \textbf{verifiable testbed for policy analysis}. 

For researchers and policy-makers, the framework proves that subjective LLM behaviors — such as deterrence breakdown, alliance formation, and bureaucratic vetoes — can be empirically measured, tracked over time, and causally audited. This provides a powerful foundation for evaluating the safety and reliability boundaries of deploying autonomous agents in high-stakes, adversarial environments.

% -------------------------------------------------------------------
\section{Limitations and Future Work}
While GeoMAS establishes a robust foundation for geopolitical simulation, the findings are subject to specific methodological boundaries that define the roadmap for future development.

First, the current experimental validation relies exclusively on a single state-of-the-art model (DeepSeek V3.2). Although this eliminates cross-model architectural noise, it introduces the risk of model-specific behavioral biases. \textbf{Future iterations} should conduct large-scale comparative studies across an ensemble of heterogeneous models to determine whether the emergent phenomena — such as Moral Washing and hegemonic unipolarity — are universal properties of current frontier models or artifacts of a specific training paradigm.

Second, the environmental mechanics governing GeoMAS, while deterministic, rely on static heuristics for resource elasticity and unrest thresholds. \textbf{Subsequent research} should explore dynamic recalibration of the physics engine, potentially employing reinforcement learning loops to automatically tune the macro-economic parameters to match real-world historical validation sets, thereby increasing the ecological validity of the simulation.

Third, while the Pandemic scenario provided crucial insights into systemic shock absorption, it represents only one dimension of exogenous perturbation. \textbf{Future work} must focus on the definition and implementation of a diverse taxonomy of new scenarios. This includes asymmetric localized shocks (e.g., targeted resource discoveries, regional natural disasters) and explicit adversarial injections (e.g., sponsored separatist insurrections or deliberate disinformation campaigns) to evaluate how differently the multi-agent system reacts to localized vs.\ global, and natural vs.\ adversarial stressors.

Finally, the cognitive architecture of the agents is currently restricted to top-down government cabinets (President, Defense Minister, Foreign Minister). A critical structural limitation of this approach is the deterministic, mathematical abstraction of the populace: internal stability is calculated by a heuristics engine rather than simulated socially. \textbf{Future developments} should integrate a \textbf{Public Opinion Agent} into each nation's hierarchy—an LLM explicitly prompted to read the state's public actions, evaluate domestic welfare, and generate narrative unrest or approval. Replacing the deterministic heuristic with a stochastic socio-political actor would complete the neuro-symbolic stack, forcing the government cabinet to negotiate not just with external adversaries, but with the evolving, autonomous ideological demands of their own citizens.

\vspace{1em}
In conclusion, GeoMAS demonstrates that the rigorous, structured simulation of LLM societies is not only possible but necessary. As artificial intelligence continues its rapid integration into the mechanisms of global statecraft, frameworks capable of safely modeling, auditing, and predicting their emergent strategic behaviors will remain an indispensable asset for ensuring international stability.
